\section{The homogeneous obstruction in degree four}\label{sec:rigidity}

For a sequence with \(X\to\infty\), Proposition~\ref{prop:ord-tightness}
shows that every weak subsequential limit of the packet measures is an
\(A_4\)-invariant probability.  Proposition~\ref{prop:ord-balls}
supplies the compact ball estimate in Equation~\eqref{eq:ord-balls}.
We first turn that estimate
into positive entropy in almost every \(A_4\)-ergodic component.
Measure classification then makes those components homogeneous.
In degree four, the proper homogeneous components that remain have two
coordinate blocks of dimension two.  We will show that each of them
carries an exterior-square lattice vector in one of three fixed planes.
The arithmetic part of the proof will exclude these three events.

\subsection{Entropy from the compact ball estimate}

Fix distinct real numbers \(\lambda_1,\ldots,\lambda_4\) with sum zero,
and set
\[
 a(t)=\diag(e^{\lambda_1t},\ldots,e^{\lambda_4t}),
 \qquad
 \lambda_a=\min_{i\ne j}|\lambda_i-\lambda_j|>0.
\]
For an \(a(\R)\)-invariant probability \(\xi\), let \(h_\xi(a(1))\)
denote the measure-theoretic entropy of right multiplication by \(a(1)\).
The argument below is the degree-four specialization of the
ordinary-ball argument in
\cite[Proposition~6.1 and Lemma~6.2]{PrimeDegreePackets}, which in turn
follows the strategy of \cite[Section~2.7.1]{ELMV2011}.
We include the argument because its statement for every dominated
invariant probability is needed when passing to ergodic components.

Here is the entropy input.  For an open identity neighborhood \(B\)
in \(\mathrm{SL}_4(\R)\), write
\[
 B^{(s,t)}=a(-s)Ba(s)\cap a(-t)Ba(t).
\]
The following is the specialization to \(X_4\) of the
Einsiedler--Lindenstrauss--Michel--Venkatesh tube criterion
\cite[Corollary~3.3 and Equation~(3.4)]{ELMV2009}.  The flow \(a(t)\)
is \(\R\)-diagonalizable, as required there.

\begin{theorem}[Entropy criterion for two-sided tubes]
\label{thm:rig-tube-criterion}
Let \(\xi_i\) and \(\xi\) be \(a(\R)\)-invariant Borel probabilities
on \(X_4\), with \(\xi_i\) converging weak-* to \(\xi\).
Suppose that \(t_i>0\) tend to infinity and that \(\eta>0\).
Assume that for every compact \(\Omega\subset X_4\) there are an open
identity neighborhood \(B\subset\mathrm{SL}_4(\R)\) and a constant
\(C_\Omega\) such that
\[
 \xi_i\bigl(xB^{(-t_i,t_i)}\bigr)
 \le C_\Omega e^{-2\eta t_i}
 \qquad(x\in\Omega,\ i\ge1).
\]
Then \(h_\xi(a(1))\ge\eta\).
\end{theorem}

In this criterion \(B\) may depend on \(\Omega\), and the
probabilities need not be ergodic under the flow.  The requirement
that the weak-* limit is a probability is part of the theorem.

\begin{proposition}[Entropy for dominated probabilities]
\label{prop:rig-dominated-entropy}
Let \(\mu\) be an \(A_4\)-invariant probability on \(X_4\) satisfying
the compact ball bound in Equation~\eqref{eq:ord-balls} for every compact set of
centers.  If \(\xi\) is an \(A_4\)-invariant probability with
\(\xi\le c\mu\) for some finite \(c\), then
\[
 h_\xi(a(1))\ge \lambda_a/3.
\]
\end{proposition}

\begin{proof}
Use the entrywise neighborhoods \(B(r)\) defined in
Equation~\eqref{eq:ord-ball-neighborhood}.  Fix \(0<b<1/4\) and put
\[
 T_t=a(t)B(b)a(-t)\cap a(-t)B(b)a(t),
 \qquad r=e^{-\lambda_a t}\quad(t\ge0).
\]
For \(h\in T_t\), the two conjugations give
\begin{equation}\label{eq:rig-tube-entries}
 |h_{ii}-1|<b,\qquad
 |h_{ij}|<b e^{-|\lambda_i-\lambda_j|t}\le br
 \quad(i\ne j).
\end{equation}
In the determinant expansion, every term except the product of the
diagonal entries contains at least two off-diagonal entries.
Since \(\det h=1\), Equation~\eqref{eq:rig-tube-entries} implies
\(\prod_i h_{ii}=1+O(r^2)\).  The diagonal matrix
\[
 d(h)=\diag\bigl(h_{11},h_{22},h_{33},
                         (h_{11}h_{22}h_{33})^{-1}\bigr)
\]
therefore belongs to a fixed compact subset
\(\mathcal D_\infty\subset A_4\), and \(d(h)^{-1}h\in B(Cr)\).
The constant \(C\) here and below may depend on \(b\) and the fixed flow.
Indeed, the first three diagonal entries of \(d(h)^{-1}h\) are one,
the fourth differs from one by \(O(r^2)\), and all its off-diagonal
entries are \(O(r)\).

The logarithmic coordinates on \(A_4\) identify
\(\mathcal D_\infty\) with a bounded subset of \(\R^3\).
A mesh of spacing comparable to \(r\), followed by a fixed enlargement
to account for matrix multiplication, gives
\begin{equation}\label{eq:rig-tube-cover}
 T_t\subset\bigcup_{j=1}^{N_t}d_jB(Cr),
 \qquad d_j\in\mathcal D_\infty,\qquad N_t\le Cr^{-3}.
\end{equation}
For a compact \(\Omega\subset X_4\), the centers \(xd_j\) lie in the
compact set \(\Omega\mathcal D_\infty\).  Domination and
Equation~\eqref{eq:ord-balls} now imply
\begin{equation}\label{eq:rig-tube-mass}
 \xi(xT_t)\le C_{\Omega,c}r^{-3}r^4
             =C_{\Omega,c}e^{-\lambda_a t}
 \qquad(x\in\Omega)
\end{equation}
for all sufficiently large \(t\).
Apply Theorem~\ref{thm:rig-tube-criterion} to the constant sequence
\(\xi_i=\xi\), with \(t_i=i\), \(B=B(b)\), and
\(\eta=\lambda_a/3\).  The identity \(B(b)^{(-t,t)}=T_t\) and
\(2\lambda_a/3<\lambda_a\) verify its estimate for large \(i\).
For the finitely many earlier \(i\), enlarge \(C_{\Omega,c}\), using
that \(\xi\) is a probability.  The criterion gives the assertion.
\end{proof}

\begin{corollary}[Positive entropy in the ergodic components]
\label{cor:rig-positive-components}
Let \(\mu\) be an \(A_4\)-invariant probability on \(X_4\) satisfying
Equation~\eqref{eq:ord-balls} on every compact set of centers.  Almost every
component in its \(A_4\)-ergodic decomposition has positive
\(a(1)\)-entropy.  In particular, this holds for every probability
limit in Proposition~\ref{prop:ord-balls}.
\end{corollary}

\begin{proof}
The group \(A_4\simeq\R^3\) is locally compact and second countable,
and acts continuously on the standard Borel space \(X_4\).
The ergodic decomposition theorem
\cite[Theorem~5.2]{GreschonigSchmidt2000}, specialized to invariant
probabilities, gives a measurable decomposition
\[
 \mu=\int_Z\nu_z\,d\tau(z)
\]
into \(A_4\)-invariant, \(A_4\)-ergodic probabilities over a factor
fixed by \(A_4\).  Put \(T(x)=xa(1)\).

We record why a restriction to the zero-entropy components is
measurable.  Choose increasing finite Borel partitions
\(\mathcal P_k\) generating the Borel \(\sigma\)-algebra of \(X_4\).
For any \(T\)-invariant probability \(\nu\), the
Kolmogorov--Sinai approximation theorem and subadditivity give
\[
 h_\nu(T)=\sup_k\inf_{N\ge1}\frac1N
 H_\nu\left(\bigvee_{j=0}^{N-1}T^{-j}\mathcal P_k\right),
\]
where \(H_\nu\) is the Shannon entropy of a finite partition.
Each finite-partition entropy is a Borel function of \(\nu\), so
\(z\mapsto h_{\nu_z}(T)\) is measurable.
For each probability under consideration, complete the standard Borel
probability space; it is then a normalized Lebesgue space.
The entropy integral formula for an invariant component factor
\cite[Theorem~9.8]{Rokhlin1967} gives
\begin{equation}\label{eq:rig-entropy-decomposition}
 h_\mu(T)=\int_Z h_{\nu_z}(T)\,d\tau(z),
\end{equation}
and the same formula after restricting and normalizing \(\tau\) on
any measurable subset of positive measure.  This formula applies
because the component factor is fixed by \(T\).  It does not require
the individual \(\nu_z\) to be \(T\)-ergodic.

If \(Z_0=\{z:h_{\nu_z}(T)=0\}\) had measure \(w>0\), then
\[
 \xi=\frac1w\int_{Z_0}\nu_z\,d\tau(z)
\]
would be an \(A_4\)-invariant probability with \(\xi\le w^{-1}\mu\).
The restricted form of Equation~\eqref{eq:rig-entropy-decomposition}
would give \(h_\xi(T)=0\), contrary to
Proposition~\ref{prop:rig-dominated-entropy}.
\end{proof}

\subsection{Finite-volume subgroups containing the diagonal group}

We next use the measure-classification theorem of
Einsiedler--Katok--Lindenstrauss
\cite[Theorem~1.3; definition in Conjecture~1.1]{EKL2006}.
We state the exact specialization that supplies the homogeneous
measures needed here.

\begin{theorem}[Positive-entropy measure classification]
\label{thm:rig-ekl-classification}
For \(n\ge3\), let \(A_n\) be the group of positive diagonal matrices
in \(\mathrm{SL}_n(\R)\).  Let \(\nu\) be an \(A_n\)-invariant and
\(A_n\)-ergodic probability on
\(\mathrm{SL}_n(\R)/\mathrm{SL}_n(\Z)\), with the left action.
If \(h_\nu(a_0(1))>0\) for some one-parameter subgroup
\(a_0:\R\to A_n\), then there are a
closed connected subgroup \(L\subset\mathrm{SL}_n(\R)\) containing
\(A_n\) and a point \(x\) such that \(Lx\) is closed and \(\nu\) is
the \(L\)-invariant probability on \(Lx\).
\end{theorem}

No ergodicity for the chosen one-parameter subgroup is assumed.
For \(n=4\), inversion
\(\mathrm{SL}_4(\Z)g\mapsto g^{-1}\mathrm{SL}_4(\Z)\)
turns right multiplication by \(a(t)\) into left multiplication by
\(a(-t)\), and entropy is unchanged upon inverting an invertible
transformation.  Thus, on our right quotient, the theorem says that
a positive-entropy \(A_4\)-ergodic probability is invariant
probability on an orbit \(\Lambda J\), where \(J\) is closed and
connected, contains \(A_4\), and the stabilizer of \(\Lambda\) is a
lattice in \(J\).

For a partition \(\mathcal I=\{I_1,\ldots,I_s\}\) of
\(\{1,2,3,4\}\) into nonempty subsets, define
\begin{equation}\label{eq:rig-block-group}
 J_{\mathcal I}
 =\left\{g\in\mathrm{SL}_4(\R):
 \begin{array}{l}
 g_{ij}=0\text{ if \(i,j\) belong to different \(I_v\)},\\
 \det(g|_{\R^{I_v}})>0\text{ for every }v
 \end{array}\right\}.
\end{equation}
This is the identity component of the corresponding block-diagonal
group.  Lindenstrauss and Weiss identified the equal-block phenomenon
for orbit closures containing a compact diagonal orbit
\cite[Theorem~1.3]{LindenstraussWeiss2001}; their result also enters the
proof of \cite[Corollary~1.4]{EKL2006}.  The next lemma gives a direct
finite-volume argument and records the rational structure of the blocks.

\begin{lemma}[Equal blocks and a rational commutant]
\label{lem:rig-equal-blocks}
Let \(J\) be a closed connected subgroup of \(\mathrm{SL}_4(\R)\)
containing \(A_4\), and let \(\Lambda\in X_4\).  Suppose that the
stabilizer \(\Gamma_\Lambda=\{j\in J:\Lambda j=\Lambda\}\) is a
lattice in \(J\).
Then \(J=J_{\mathcal I}\) for a partition \(\mathcal I\) whose blocks
all have the same size.

More precisely, put \(W_\Q=\Q\Lambda\) and
\(W_\R=W_\Q\otimes_\Q\R=\R^4\).  There is a totally real number field
\(E\subset\End_\Q(W_\Q)\), of degree \(s=|\mathcal I|\), whose scalar
extension to \(\R\) is exactly the commutant of \(J\) in
\(\End_\R(W_\R)\).  The coordinate blocks \(\R^{I_v}\) are the real
eigenspaces of \(E\) on \(W_\R\).
\end{lemma}

\begin{proof}
Let \(\mathfrak j\) be the Lie algebra of \(J\).
It contains the trace-zero diagonal algebra and is invariant under
diagonal conjugation.  The other weight spaces in
\(\mathfrak{sl}_4(\R)\) are the distinct one-dimensional spaces
\(\R E_{ij}\), \(i\ne j\), where \(E_{ij}\) is a matrix unit.
An invariant subspace is a sum of these weight spaces, so
\(\mathfrak j\) consists of the diagonal algebra and a selection of
the root spaces.

The group \(J\) is unimodular.
One way to see this standard fact is to use the quotient
criterion for Haar measures: existence of an invariant quotient
measure and triviality of the modular character of the discrete
group \(\Gamma_\Lambda\) imply that the modular character of \(J\)
is trivial on \(\Gamma_\Lambda\).  It therefore descends to
\(\Gamma_\Lambda\backslash J\).  Its pushforward of invariant
probability would be invariant under translations by its image.
A nontrivial positive dilation admits no invariant probability on
\(\R_{>0}\), as its disjoint dilates of a fundamental interval all
have the same mass.  Hence the modular character is trivial on \(J\).

Draw a directed edge \(i\to j\) when \(E_{ij}\in\mathfrak j\).
The trace of the adjoint action of
\(\diag(u_1,\ldots,u_4)\), with \(\sum_i u_i=0\), on
\(\mathfrak j\) is the sum of the selected roots \(u_i-u_j\).
Unimodularity makes this sum zero.  Its coefficient at \(u_i\) is
the outgoing degree of \(i\) minus its incoming degree.  These
coefficients sum to zero, and a linear form vanishing on
\(\sum_i u_i=0\) has all coefficients equal.  Thus every vertex
has equal outgoing and incoming degrees.

Every edge now lies within a strongly connected component.
Otherwise the finite directed acyclic graph of components would
have a source with an outgoing edge; summing degree balance over
that source would give a positive outgoing degree and zero incoming
degree.  For distinct indices, the bracket identity
\([E_{ij},E_{jk}]=E_{ik}\) makes the edge relation transitive.
Each strongly connected component is consequently a complete
directed graph.  These components form a partition \(\mathcal I\),
and \(\mathfrak j\) is the Lie algebra of \(J_{\mathcal I}\).
Connected Lie subgroups with the same Lie algebra agree, giving
\(J=J_{\mathcal I}\).

We use \(\Gamma_\Lambda\) to recover a rational structure on the
block scalars.  Its elements preserve \(\Lambda\), hence act
rationally on \(W_\Q\).  Let
\[
 \mathcal C_\Q
 =\{M\in\End_\Q(W_\Q):M\gamma=\gamma M
                         \text{ for every }\gamma\in\Gamma_\Lambda\}.
\]
In a lattice basis the commutation equations have integer
coefficients.  A finite subcollection of these linear equations
defines the same solution space, by finite dimensionality.
It follows that \(\mathcal C_\Q\otimes_\Q\R\) is the real
commutant of \(\Gamma_\Lambda\).

We claim that this is also the real commutant of \(J\).
Let \(M\) be a real endomorphism commuting with \(\Gamma_\Lambda\).
The continuous map
\[
 \Gamma_\Lambda\backslash J\longrightarrow\End_\R(W_\R),
 \qquad \Gamma_\Lambda j\longmapsto j^{-1}Mj
\]
is well-defined.  Its pushforward of invariant probability is
invariant under conjugation by \(A_4\).  Each off-diagonal
coordinate is multiplied by a nontrivial positive dilation under
some element of \(A_4\).  A probability on \(\R\) invariant under
such a dilation is supported at zero: the annuli
\(\{x:r^k\le |x|<r^{k+1}\}\), \(k\in\Z\), for a fixed \(r>1\),
are disjoint dilates covering \(\R\setminus\{0\}\).
Thus every off-diagonal coordinate of \(j^{-1}Mj\) is zero almost
everywhere.  Invariant probability on the transitive space
\(\Gamma_\Lambda\backslash J\) has full support; continuity makes
\(j^{-1}Mj\) diagonal for every \(j\in J\).
In particular \(M\) is diagonal.  Conjugating by
\(1+tE_{ij}\) for \(i,j\) in the same block forces its \(i\)-th
and \(j\)-th diagonal entries to agree.  Therefore \(M\) is scalar
on each block, which is exactly the condition that it commute with
\(J\).  The reverse inclusion is immediate, proving the claim.

It follows that \(\mathcal C_\Q\otimes_\Q\R\simeq\R^s\).
In particular, \(\mathcal C_\Q\) is commutative and has no
nonzero nilpotents.  The structure theorem for finite reduced
commutative algebras over \(\Q\) expresses it as a product of
number fields, all totally real because its real scalar extension
has no complex factor.

This product has only one factor.  If
\(e\in\mathcal C_\Q\) were an idempotent other than \(0\) and \(1\),
then \(W=e(W_\R)\) would be a nonempty proper union of coordinate
blocks, rational relative to \(\Lambda\).  Every
\(\gamma\in\Gamma_\Lambda\) preserves the full lattice
\(\Lambda\cap W\) in \(W\), so
\(\lvert\det(\gamma|_W)\rvert=1\).
The continuous character
\[
 \chi:J\longrightarrow\R_{>0},\qquad
 \chi(j)=\lvert\det(j|_W)\rvert
\]
is therefore trivial on \(\Gamma_\Lambda\).
It is surjective: if \(k=\dim W\), a diagonal matrix may scale the
coordinates in \(W\) by \(e^{(4-k)t}\) and those in its complement
by \(e^{-kt}\).  Its total determinant is one and its determinant
on \(W\) is \(e^{k(4-k)t}\).
The character \(\chi\) would consequently push the invariant
probability on \(\Gamma_\Lambda\backslash J\) to a probability on
\(\R_{>0}\) invariant under all multiplicative translations.
The disjoint-interval argument above rules this out.
Thus \(\mathcal C_\Q\) is a field, which we denote by \(E\).

Since \(E\otimes_\Q\R\simeq\R^s\), its degree is \(s\).
The space \(W_\Q\) is a vector space over \(E\), and
\[
 W_\Q\otimes_\Q\R
 \simeq\bigoplus_{\tau:E\hookrightarrow\R}
                 W_\Q\otimes_{E,\tau}\R.
\]
Each summand has dimension \(\dim_E W_\Q=4/s\).
These summands are the eigenspaces of the block-scalar algebra,
namely the coordinate blocks.  The blocks therefore have equal size.
\end{proof}

For a positive-entropy homogeneous probability the lemma leaves
\(J=\mathrm{SL}_4(\R)\) or two blocks of size two.
Indeed, the only other equal-block possibility is four singleton
blocks, giving \(J=A_4\).  A lattice in \(A_4\simeq\R^3\) has
compact quotient, and translations on this torus have zero entropy:
for a flat metric they are isometries, so the iterated Bowen metrics
equal the original metric and the topological entropy is zero;
the variational principle bounds measure entropy by topological entropy.

For a possible two-block limit orbit \(\Lambda J\), the quadratic field
\(E\subset\End_\Q(\Q\Lambda)\) has real scalar extension equal to
the commutant of \(J\) in \(\End_\R(\R^4)\).
The construction supplies no embedding of \(E\) into the packet fields
\(K_i\), so their primitivity does not by itself exclude this orbit.
We will exclude it by proving that packet limits give zero mass to the
exterior obstruction below.

\subsection{The three exterior obstruction events}

Let \(e_1,\ldots,e_4\) be the coordinate basis of \(\R^4\).
For \(w\in\bigwedge^2\R^4\), write \(w_{ij}\), \(i<j\), for its
coefficients in the basis \(e_i\wedge e_j\), and define
\begin{equation}\label{eq:rig-Q}
 Q(w)=w_{12}w_{34}-w_{13}w_{24}+w_{14}w_{23}.
\end{equation}
The identity
\[
 w\wedge w=2Q(w)\,e_1\wedge e_2\wedge e_3\wedge e_4
\]
shows that
\(Q((\bigwedge^2g)w)=\det(g)Q(w)\) for \(g\in\mathrm{GL}_4(\R)\).
Expressing a unimodular lattice in a determinant-one basis gives
\begin{equation}\label{eq:rig-Q-integral}
 Q\bigl(\bigwedge\nolimits^2\Lambda\bigr)\subset\Z
 \qquad(\Lambda\in X_4).
\end{equation}
This also covers the sign translates in the packet definition:
a coordinate reflection of determinant \(-1\) changes only the
sign of \(Q\).

Let
\[
 \mathcal B=\{12|34,\ 13|24,\ 14|23\}.
\]
For \(b=ij|kl\in\mathcal B\), with the indices within each pair
increasing, put
\begin{equation}\label{eq:rig-planes}
 V_b=\operatorname{span}_\R\{e_i\wedge e_j,e_k\wedge e_l\}.
\end{equation}
The restriction of \(Q\) to \(V_b\) is, up to sign, the product of
its two displayed coordinates.  Define
\begin{equation}\label{eq:rig-events}
 \mathcal E_b
 =\left\{\Lambda\in X_4:
   \text{there exists }w\in\bigwedge\nolimits^2\Lambda\cap V_b
   \text{ with }Q(w)\ne0\right\}.
\end{equation}
Each \(\mathcal E_b\) is \(A_4\)-invariant and Borel.
For Borel measurability, take a countable cover of \(X_4\) by local
charts admitting continuous lattice bases.  In a chart, the six
integer coefficients of an exterior vector label it continuously.
For each label, membership in \(V_b\) is a set of continuous
coordinate equalities, and \(Q(w)\ne0\) is an open condition.
Taking the countable union over labels and charts proves the claim.
Invariance follows because \(A_4\) preserves \(V_b\) and \(Q\).

\begin{proposition}[The quartic homogeneous obstruction]
\label{prop:rig-obstruction}
Let \(\nu\) be an \(A_4\)-invariant, \(A_4\)-ergodic probability
on \(X_4\).  Suppose that \(h_\nu(a_0(1))>0\) for some
one-parameter subgroup \(a_0:\R\to A_4\).
Then either \(\nu=m_4\), or there exists
\(b\in\mathcal B\) such that \(\nu(\mathcal E_b)=1\).
In the latter case every lattice \(\Lambda'\) in the homogeneous orbit carrying
\(\nu\) has a vector in \(\bigwedge^2\Lambda'\cap V_b\) with both
coordinates nonzero and with nonzero integer \(Q\)-value.
\end{proposition}

\begin{proof}
Theorem~\ref{thm:rig-ekl-classification}, in the right-quotient
form explained above, represents \(\nu\) as invariant probability
on \(\Lambda J\) for a connected closed subgroup \(J\) containing
\(A_4\).  By Lemma~\ref{lem:rig-equal-blocks} and the
positive-entropy assumption, either \(J=\mathrm{SL}_4(\R)\), which
gives \(\nu=m_4\), or \(J\) has two coordinate blocks of size two.

In the second case write the blocks as \(\{i,j\}\) and \(\{k,l\}\),
and put \(b=ij|kl\).  The commutant field \(E\) in
Lemma~\ref{lem:rig-equal-blocks} is quadratic.  Choose a nonzero
\(\alpha\in E\) with \(\Tr_{E/\Q}(\alpha)=0\).
Its two scalar values on the blocks are \(c\) and \(-c\), with
\(c\ne0\).  The rational endomorphism \(\alpha\) of
\(W_\Q=\Q\Lambda\) induces a rational endomorphism
\[
 D_\alpha:\bigwedge\nolimits^2 W_\Q\longrightarrow
                    \bigwedge\nolimits^2 W_\Q,\qquad
 D_\alpha(v\wedge w)=\alpha(v)\wedge w+v\wedge\alpha(w).
\]
On the exterior square of the first block it acts by \(2c\), and
on that of the second by \(-2c\).  It vanishes on the cross-block
summand.  Its real image is therefore exactly
\[
 \bigwedge\nolimits^2\R^{\{i,j\}}
 \oplus \bigwedge\nolimits^2\R^{\{k,l\}}=V_b.
\]
Because \(D_\alpha\) is rational, \(V_b\) is a rational subspace
relative to \(\bigwedge^2 W_\Q\).

It follows that \(V_b\cap\bigwedge^2\Lambda\) is a rank-two lattice
in \(V_b\).  It contains a vector off both coordinate axes: take
two independent lattice vectors; if neither is already off the
axes, their sum is.  For this vector \(Q(w)\ne0\), and
Equation~\eqref{eq:rig-Q-integral} makes its value a nonzero integer.
Every element of \(J\) preserves each of the two exterior lines,
multiplying its coefficient by the nonzero determinant of the
corresponding block.  The translated vector therefore has both
coordinates nonzero at every lattice of \(\Lambda J\).
This proves \(\nu(\mathcal E_b)=1\).
\end{proof}

Corollary~\ref{cor:rig-positive-components} and
Proposition~\ref{prop:rig-obstruction} reduce identification of a
packet probability limit \(\mu\) to proving
\(\mu(\mathcal E_b)=0\) for all three \(b\in\mathcal B\).
Equation~\eqref{eq:rig-Q-integral} lets us test those events one
nonzero integer \(Q\)-value at a time.  The next sections establish
the exterior-vector estimate used for this exclusion.
